\subsection{Spinal Load Series and Sign Conventions}
\label{sec:spn-5}

This section documents the load container and the tension-positive bridge of
the segmental spine (issue \#12747, epic \#12735, ADR-0055 item 5). The code is
\texttt{src/shared/python/spine/loads.py}; it performs no dynamics, only
bookkeeping, frame rotation and wrench transport of loads produced elsewhere.

\paragraph{Observables.}
Three observables are carried and never conflated: \texttt{GENERALIZED\_EFFORT}
(inverse-dynamics generalized force, all structures lumped),
\texttt{JOINT\_REACTION} (OpenSim \texttt{JointReaction} with static-optimization
muscle forces applied) and \texttt{PASSIVE\_ELEMENT} (bushing, ligament and disc
stiffness elements only).

\paragraph{Frames and Signs.}
For every observable the wrench $(\mathbf F,\mathbf M)$ is that exerted
\emph{on the superior vertebra by the inferior vertebra}, expressed in the
inferior vertebra ISB frame ($X$ anterior, $Y$ cranial, $Z$ right), about the
mid-disc joint centre. Units are N and N\,m.
\begin{align}
  \text{compression} &= +F_y, &
  \text{AP shear} &= +F_x\ (\text{anterior}), &
  \text{ML shear} &= +F_z\ (\text{right}).
\end{align}
A superior body of mass $m$ at rest is supported by $F_y=+mg$ (compression
positive). Tension-positive consumers
(\texttt{body\_part\_viz/axial\_loads.py}) receive
$T=-F_y$.

\paragraph{Wrench Transport.}
Given source-frame samples $(\mathbf F_s,\mathbf M_s)$ applied at a reference
point displaced by $\mathbf r$ from the joint centre ($\mathbf r$ expressed in
the inferior ISB frame) and a rotation $R$ with $\mathbf v_{\mathrm{ISB}}=R\,\mathbf v_s$,
\begin{equation}
  \mathbf F = R\mathbf F_s,\qquad
  \mathbf M = R\mathbf M_s + \mathbf r\times\mathbf F ,
\end{equation}
evaluated through \texttt{move\_wrench\_point}. Worked check: $\mathbf F=(0,100,0)$\,N
at $\mathbf r=(0.05,0,0)$\,m gives $\mathbf M=(0,0,5)$\,N\,m.

\paragraph{Normalisation.}
\texttt{in\_body\_weights} divides forces by $mg$ with $g=9.80665$\,m/s$^2$
(\texttt{GRAVITY\_FLOAT}) and moments by $mg$ (units BW\,m). The series then
reports \texttt{force\_unit="BW"} and records the mass in provenance; it cannot
feed the newton-valued axial bridge.

\paragraph{Peaks.}
\texttt{peak(q)} returns, per level, the signed value of largest magnitude and
its time; ties resolve to the earliest sample.

\paragraph{Provenance.}
Every series requires non-empty string values for \texttt{model\_id},
\texttt{model\_sha256}, \texttt{solver}, \texttt{cost\_exponent} (\texttt{n/a}
when unused), \texttt{distribution\_source} and \texttt{external\_loads}.
Missing or blank keys raise \texttt{ValueError}.

\paragraph{Preconditions.}
Finite, strictly increasing 1-D time; distinct levels; force and moment of
shape $(n,L,3)$ and finite; rotations proper ($R^\top R=I$, $\det R=1$,
tolerance $10^{-6}$). Arrays are stored read-only.

\paragraph{Axial Bridge.}
\texttt{SpinalAxialLoadProvider(series, time\_s)} satisfies the
\texttt{AxialLoadProvider} protocol at one sample time (no interpolation; the
time must match a sample within $10^{-9}$\,s) and returns tension $=-$compression
per level name.

\paragraph{Limitations.}
No load values are computed or qualified here; accuracy depends on the
producing analysis, and a successful container build is not physical
acceptance. Cervical levels are out of scope.

\paragraph{Reproduction.}
\begin{verbatim}
python -m pytest tests/shared/python/spine/test_loads.py -q -n 0
\end{verbatim}
